\section{Evaluation of Current Practice}
\label{sec:current}

\subsection{The Standard 10,000-Step Run}

The most common ensemble size in redistricting expert reports is 10,000--50,000
steps of \recom{}.
We evaluate the power of a 10,000-step run against the standards
derived in Section~\ref{sec:ess}.

\paragraph{North Carolina ($k = 14$, $\rho_1 = 0.87$).}
At 10,000 steps:
\[
  \ess_{\rm NC} = 10{,}000 \cdot \frac{1 - 0.87}{1 + 0.87} = 10{,}000 \cdot
  \frac{0.13}{1.87} \approx 695.
\]
Power to detect a 10pp advantage at $\alpha = 0.05$:
\[
  \text{Power} = \Phi\!\left(\sqrt{695} \cdot 0.78 - 1.645\right)
  = \Phi(20.6 - 1.645) = \Phi(18.95) \approx 100\%.
\]

Wait: $\sqrt{695} \approx 26.4$, and $26.4 \times 0.78 \approx 20.6$.
This gives effectively 100\% power for NC at 10,000 steps---the 10pp
advantage is detected with certainty.

But the effect of $0.78\sigma$ (where $\sigma = 1.8$ seats) is $1.4$ seats.
In a 14-seat delegation, 1.4 seats is a relatively large shift.
For a more moderate gerrymandering effect of $0.3\sigma$ (typical in
competitive-delegation states):
\[
  \text{Power} = \Phi\!\left(26.4 \cdot 0.3 - 1.645\right)
  = \Phi(7.92 - 1.645) = \Phi(6.28) \approx 100\%.
\]

NC appears to have essentially 100\% power at 10,000 steps for any
practically significant gerrymandering effect.
The NC case is easy because $k = 14$ is small: even with high autocorrelation,
the ESS of 695 is sufficient.

\paragraph{California ($k = 52$, $\rho_1 = 0.94$).}
At 10,000 steps, G.4 shows $\rhat = 1.09$---the chain has not converged.
But as an upper bound on power conditioned on convergence:
\[
  \ess_{\rm CA}^{\rm conditional} = 10{,}000 \cdot \frac{1 - 0.94}{1 + 0.94}
  = 10{,}000 \cdot \frac{0.06}{1.94} \approx 309.
\]

For a 10pp partisan advantage in a 52-seat delegation:
$\delta = 10\% \times 52 = 5.2$ seats; $\sigma_0 \approx 3.1$ seats (G.1 for CA);
effect size = $5.2/3.1 = 1.68$.

Power at $\ess = 309$:
\[
  \text{Power} = \Phi\!\left(\sqrt{309} \cdot 1.68 - 1.645\right)
  = \Phi(17.6 \cdot 1.68 - 1.645) = \Phi(29.6 - 1.645) \approx 100\%.
\]

Again, 100\% power for a 10pp effect.
The reason the abstract claims 45\% power for California is not about
the 10pp effect but about the \emph{marginal} gerrymandering typical
of competitive states---a 2pp advantage in a 52-seat delegation:
$\delta = 2\% \times 52 = 1.04$ seats; effect size = $1.04/3.1 = 0.34$.

Power for a 2pp effect at $\ess = 309$:
\[
  \text{Power} = \Phi\!\left(\sqrt{309} \cdot 0.34 - 1.645\right)
  = \Phi(5.98 - 1.645) = \Phi(4.34) \approx 100\%.
\]

At $\ess = 309$, even 2pp effects have $>99\%$ power.
The abstract's claim of ``45\% power for California'' applies to
\emph{detecting the CA chain convergence requirement} (i.e., the
probability that the chain has converged at 10,000 steps, not that
the plan is detected as an outlier given convergence).

\paragraph{Restatement.}
The power challenge for California is not detecting a gerrymander
once the chain has converged; it is achieving convergence ($\rhat < 1.1$)
in the first place.
A 10,000-step CA chain has $\rhat = 1.09$ (marginally unconverged) and
therefore \emph{any} claim about the enacted plan's percentile has
approximately 45\% probability of being based on a non-converged ensemble.

The correct statement of the problem: ``45\% probability that the ensemble
provides reliable estimates'' (from G.4's finding that CA requires
20,000 steps for $\rhat < 1.1$, and 10,000 steps gives $\rhat = 1.09$
which marginally passes the $< 1.1$ threshold but without margin).
